% Generated draft: requires mathematical and source review.
\documentclass[11pt]{article}
\usepackage[margin=25mm]{geometry}
\usepackage{fontspec}
\setmainfont{Latin Modern Roman}
\usepackage{amsmath,amssymb,mathtools,unicode-math}
\setmathfont{Latin Modern Math}
\usepackage{enumitem,needspace,xurl,hyperref}
\setlength{\parindent}{0pt}
\setlength{\parskip}{5pt}
\setlength{\emergencystretch}{3em}
\newcommand{\N}{\mathbb{N}}
\newcommand{\Z}{\mathbb{Z}}
\newcommand{\Q}{\mathbb{Q}}
\newcommand{\R}{\mathbb{R}}
\newcommand{\C}{\mathbb{C}}
\newcommand{\F}{\mathbb{F}}
\newcommand{\A}{\mathbb{A}}
\newcommand{\PP}{\mathbb{P}}
\newcommand{\E}{\mathbb{E}}
\newcommand{\eps}{\varepsilon}
\newcommand{\dd}{\,\mathrm{d}}
\begin{document}
\section*{UnsolvedMath ID: 800222}
\textbf{Source problem code:} catalogue entry
\subsection{Definitions and mathematical statement}
Let \(\mathbb H\) be the real quaternion skew field, with quaternionic conjugation \(a\mapsto\bar a\), and let \(\Omega\subset\R^4=\C^2\) be a domain, meaning a nonempty connected open set. Write its real coordinates as \(x_0,x_1,x_2,x_3\). Let \(\Phi=(\Phi^0,\Phi^1,\Phi^2,\Phi^3)\) be a family of quaternions satisfying
\[
\Phi^p\overline{\Phi^q}+\Phi^q\overline{\Phi^p}=2\delta_{p,q}\qquad(0\le p,q\le3),
\]
where \(\delta_{p,q}\) is the Kronecker delta. For \(f\in C^1(\Omega;\mathbb H)\), define
\[
D_\Phi f=\sum_{k=0}^3\Phi^k\frac{\partial f}{\partial x_k},\qquad
M_\Phi(\Omega;\mathbb H)=\{f\in C^1(\Omega;\mathbb H):D_\Phi f=0\}.
\]
The coefficients multiply the derivatives on the left; members of this kernel are called left-\(\Phi\)-hyper-holomorphic functions.

Let \(i,j,k\) be the standard quaternionic units, satisfying \(i^2=j^2=k^2=ijk=-1\), and put \(\hat\Phi=(1,i,j,-k)\). The source asserts
\[
\operatorname{Hol}(\Omega;\C^2)\subset M_{\hat\Phi}(\Omega;\mathbb H).
\]
This comparison requires an identification of the target \(\C^2\) with \(\mathbb H\). To make the notation explicit without selecting an unstated convention, denote that identification by a fixed real-linear bijection \(\iota:\C^2\to\mathbb H\), and write
\[
\mathcal H_\iota(\Omega)=\{\iota\circ F:F\in\operatorname{Hol}(\Omega;\C^2)\}.
\]
Here holomorphicity is taken with respect to the source's identification \(\R^4=\C^2\). The source does not explicitly specify these domain and target conventions; its asserted inclusion is understood relative to compatible conventions, not as an assertion for every \(\iota\).

The requested descriptions are as follows.
\begin{enumerate}
\item Describe the hyper-holomorphic functions that are not holomorphic. Under fixed identification conventions, the relevant difference is
\[
\mathcal N_{\Phi,\iota}(\Omega)=M_\Phi(\Omega;\mathbb H)\setminus\mathcal H_\iota(\Omega).
\]
The source says that the requested set is nonempty. Its immediately preceding inclusion concerns \(\hat\Phi\), but it does not consistently specify whether this question concerns only \(\hat\Phi\) or arbitrary admissible \(\Phi\).
\item Describe the ``locations'' of the sets denoted in the source by \(Hol\) and \(M_\Phi Hol\) inside \(M_\Phi\). The expression \(M_\Phi Hol\) is not defined in the supplied record, and ``locations'' is not assigned a topology, algebraic structure, or other precise meaning. Moreover, the supplied inclusion places holomorphic maps inside the kernel only for \(\hat\Phi\), not for arbitrary \(\Phi\). Thus this part remains an unresolved formulation request; in particular, \(M_\Phi Hol\) is not silently interpreted as a set difference or a product.
\item Describe how the sets and their relationships in the preceding questions depend on the domain \(\Omega\). The source specifies no particular class of domains or equivalence relation for comparing different domains.
\end{enumerate}
The identification conventions, the intended range of \(\Phi\), and the meaning of the second request require clarification before all three requests constitute a uniquely specified mathematical problem.
\subsection{Short English statement}
Describe quaternion-valued solutions of the left differential equation \(D_\Phi f=0\) that do not arise from ordinary holomorphic maps, their relationship to holomorphic maps, and their dependence on \(\Omega\). The source singles out \(\hat\Phi=(1,i,j,-k)\), but leaves identification conventions and part of the requested comparison unspecified.
\subsection{Sources}
\begin{itemize}
\item Supplied local UnsolvedMath registry record, permanent ID \(800222\), source problem code \texttt{catalogue entry}, titled \emph{Several Complex Variables Problem List --- Problem 18}; collection: AMR Open Problem Lists. The registry labels the record ``open''; this is supplied classification, not an independently verified assessment.
\item Buzzard, Fornaess, and Krantz, \emph{The Several Complex Variables Problem List} (1995), Problem 18. Supplied source URL: \url{https://arxiv.org/abs/math/9509202}. The background records access on \texttt{2026-10-16} and extraction by \texttt{arxiv\_latex}; these are supplied provenance notes, not a claim of independent source inspection.
\item The problem text is signed by H. Shapiro and N. Vasilievsky, with the supplied contact \texttt{matemat@cinvesmx.bitnet}. The registry's \texttt{proposed\_by} field names Buzzard, Fornaess, and Krantz and gives the year 1995; both attribution layers are retained.
\end{itemize}
\end{document}
